% !TeX root = ../main.tex
\chapter{TỔNG QUAN VỀ XÁC THỰC TẬP TRUNG VÀ TRIỂN KHAI TỰ ĐỘNG}
\label{ch:tongquan}

\section{Xác thực, định danh và quản lý truy cập}
\label{sec:iam}

\subsection{Khái niệm xác thực và ủy quyền}

Xác thực (Authentication) là quá trình kiểm tra danh tính của một chủ
thể (người dùng hoặc dịch vụ) nhằm trả lời câu hỏi ``bạn là ai?''. Ba yếu tố
xác thực phổ biến gồm:
\begin{itemize}
  \item \textbf{Yếu tố kiến thức} (something you know): mật khẩu, mã PIN.
  \item \textbf{Yếu tố sở hữu} (something you have): điện thoại, token, khóa
  bảo mật.
  \item \textbf{Yếu tố sinh trắc} (something you are): vân tay, khuôn mặt.
\end{itemize}

Khi kết hợp từ hai yếu tố trở lên, ta có xác thực đa yếu tố (MFA).
Khác với xác thực, ủy quyền (Authorization) trả lời câu hỏi ``bạn được
phép làm gì?'' và thường được mô hình hóa bằng vai trò (RBAC) hoặc chính sách
chi tiết (ABAC). Xác thực và ủy quyền cùng thuộc lĩnh vực quản lý định danh và truy cập (IAM).

\subsection{Đăng nhập một lần và định danh liên kết}

Đăng nhập một lần (SSO) cho phép người dùng xác thực một lần và truy
cập nhiều ứng dụng mà không phải đăng nhập lại. Trong mô hình
định danh liên kết (Federated Identity), một nhà cung cấp định danh (IdP) chịu trách nhiệm xác thực, còn các ứng dụng -- gọi là
bên tin cậy (Relying Party -- RP) -- tin tưởng vào kết quả xác thực của
IdP thông qua các giao thức chuẩn như OAuth 2.0, OpenID Connect hoặc SAML.
Mô hình này cho phép tách biệt trách nhiệm xác thực khỏi logic nghiệp vụ của
từng ứng dụng.

% TODO: Có thể thêm hình minh họa mô hình SSO (3 vai trò: User, IdP, RP).

\subsection{Các mối đe dọa phổ biến đối với hệ thống xác thực}
\label{subsec:threats}

Do nắm giữ ``chìa khóa'' truy cập vào mọi ứng dụng, hệ thống xác thực tập trung
là mục tiêu hấp dẫn của kẻ tấn công. Một số mối đe dọa điển hình gồm:

\begin{itemize}
  \item \textit{Phishing}: tạo trang đăng nhập giả mạo để lừa người dùng nhập
  tài khoản và mật khẩu.
  \item \textit{Brute force} và \textit{credential stuffing}: thử liên tục
  nhiều mật khẩu, hoặc dùng thông tin đăng nhập rò rỉ từ các vụ vi phạm dữ liệu
  khác.
  \item \textit{Session hijacking}: đánh cắp phiên đã xác thực (cookie, token)
  để mạo danh người dùng mà không cần mật khẩu.
  \item \textit{Man-in-the-middle}: chặn bắt thông tin trên đường truyền nếu
  không được mã hóa.
  \item \textit{Token replay}: đánh cắp và sử dụng lại token còn hiệu lực.
  \item \textit{Privilege escalation}: lợi dụng cấu hình sai để nâng cao quyền.
\end{itemize}

Các biện pháp giảm thiểu tương ứng bao gồm: bắt buộc HTTPS/TLS, xác thực đa
yếu tố, giới hạn số lần đăng nhập (rate limiting), đặt thời hạn token ngắn, sử
dụng PKCE, ghi log kiểm toán và áp dụng nguyên tắc đặc quyền tối thiểu. Trong
phạm vi chuyên đề, các biện pháp này được thể hiện xuyên suốt ở phần thiết kế
và triển khai.

\subsection{Lưu trữ thông tin xác thực an toàn}
\label{subsec:hashing}

Một sai lầm nghiêm trọng là lưu mật khẩu dưới dạng thô (plaintext). Thay vào
đó, mật khẩu phải được \textit{băm} (hash) kèm một giá trị ngẫu nhiên gọi là
\textit{muối} (salt) bằng các thuật toán chuyên dụng dành cho mật khẩu như
bcrypt, scrypt hoặc Argon2. Các thuật toán này được thiết kế để chậm một cách
có chủ đích, khiến việc dò ngược trở nên tốn kém.

Keycloak hỗ trợ nhiều thuật toán băm và cho phép cấu hình chính sách mật khẩu
như độ dài tối thiểu, yêu cầu ký tự, thời hạn và số lần đăng nhập sai. Việc sử
dụng một hệ thống đã được kiểm chứng thay vì tự cài đặt cơ chế mã hóa là
nguyên tắc quan trọng nhằm tránh các lỗ hổng phát sinh.

\section{Giao thức OAuth 2.0 và OpenID Connect}
\label{sec:oidc}

\subsection{OAuth 2.0}

OAuth 2.0~\cite{rfc6749} là một khung ủy quyền cho phép ứng dụng bên thứ ba
truy cập tài nguyên có kiểm soát thay mặt người dùng mà không cần biết mật
khẩu của họ. Các vai trò chính gồm: \textit{Resource Owner} (người dùng),
\textit{Client} (ứng dụng), \textit{Authorization Server} (máy chủ ủy quyền) và
\textit{Resource Server}. OAuth 2.0 định nghĩa nhiều luồng cấp quyền
(\textit{grant types}), trong đó luồng \textit{Authorization Code} là luồng phù
hợp nhất cho ứng dụng web.

\subsection{OpenID Connect}

OpenID Connect (OIDC)~\cite{oidccore} được xây dựng trên nền OAuth 2.0, bổ
sung lớp định danh thông qua \textit{ID Token} -- một JSON Web Token
(JWT)~\cite{rfc7519} chứa các thông tin về người dùng. OIDC chuẩn hóa các
điểm cuối (endpoint) như: \texttt{/authorize}, \texttt{/token},
\texttt{/userinfo} và \texttt{/.well-known/openid-configuration}, đồng thời
định nghĩa các \textit{scope} như \texttt{openid}, \texttt{profile},
\texttt{email}.

\subsection{Luồng Authorization Code và PKCE}

Luồng \textit{Authorization Code} hoạt động theo các bước tổng quát sau:

\begin{enumerate}
  \item Người dùng truy cập ứng dụng (RP) và được chuyển hướng tới IdP.
  \item IdP xác thực người dùng, sau đó chuyển hướng trở lại RP kèm theo một
  \textit{authorization code} ngắn hạn.
  \item RP trao đổi \textit{authorization code} này với IdP để lấy
  \textit{access token} và \textit{ID Token}.
  \item RP sử dụng token để truy cập tài nguyên hoặc xác lập phiên cho người
  dùng.
\end{enumerate}

Để tăng cường bảo mật cho các ứng dụng công khai, phần mở rộng
PKCE~\cite{rfc7636} (Proof Key for Code Exchange) được sử dụng: RP sinh một
cặp giá trị \texttt{code\_verifier} và \texttt{code\_challenge}, gửi
\texttt{code\_challenge} khi bắt đầu luồng và gửi \texttt{code\_verifier} khi
đổi token, nhằm chống lại tấn công đánh chặn authorization code.

% TODO: Chèn biểu đồ tuần tự (sequence diagram) của luồng Authorization Code + PKCE.

\subsection{Cấu trúc JWT và xác thực token}
\label{subsec:jwt}

JSON Web Token gồm ba phần được mã hóa Base64URL và ngăn cách bởi dấu chấm:
\textit{header} (thuật toán và loại token), \textit{payload} (các khai báo --
claims) và \textit{signature} (chữ ký số). Một số claim chuẩn quan trọng gồm:
\texttt{iss} (nhà phát hành), \texttt{sub} (chủ thể), \texttt{aud} (đối tượng
nhận), \texttt{exp} (thời điểm hết hạn), \texttt{iat} (thời điểm phát hành),
\texttt{nbf} (thời điểm bắt đầu hiệu lực) và \texttt{jti} (định danh token).

Khi nhận token, ứng dụng cần xác thực chữ ký bằng khóa công khai lấy từ điểm
cuối JWKS của IdP, đồng thời kiểm tra các claim về nhà phát hành, đối tượng và
thời hạn. Việc kiểm tra đầy đủ các bước này giúp chống lại token giả mạo hoặc
token đã hết hiệu lực.

\subsection{Các loại token và vòng đời}
\label{subsec:tokenlife}

Trong OIDC/OAuth 2.0 tồn tại ba loại token chính: \textit{access token} (dùng
để truy cập tài nguyên, thời hạn ngắn), \textit{refresh token} (dùng để xin
access token mới khi hết hạn, thời hạn dài hơn) và \textit{ID Token} (mang
thông tin định danh người dùng). Việc đặt thời hạn ngắn cho access token giúp
giảm thiệt hại nếu token bị đánh cắp, trong khi cơ chế refresh token giúp duy
trì trải nghiệm người dùng mà không phải đăng nhập lại.

\subsection{So sánh các luồng xác thực OIDC}
\label{subsec:flows}

Bảng~\ref{tab:oidcflows} tóm tắt ưu nhược điểm của các luồng phổ biến và
trường hợp sử dụng phù hợp.

\begin{table}[htbp]
\centering
\caption{So sánh các luồng xác thực OIDC}
\label{tab:oidcflows}
\small
\begin{tabular}{p{3cm} p{4.6cm} p{5.4cm}}
\toprule
\textbf{Luồng} & \textbf{Ưu điểm} & \textbf{Hạn chế / phù hợp}\\
\midrule
Authorization Code + PKCE & An toàn nhất, code trao đổi ở kênh sau & Ứng dụng web và di động\\
Implicit & Đơn giản, ít bước & Không an toàn; đã bị khuyến nghị loại bỏ\\
Hybrid & Kết hợp code và token & Phức tạp; ít dùng\\
Client Credentials & Không cần người dùng & Giao tiếp giữa các dịch vụ\\
\bottomrule
\end{tabular}
\end{table}

Trên cơ sở so sánh, chuyên đề lựa chọn luồng \textit{Authorization Code} kèm
PKCE cho ứng dụng demo, đây là khuyến nghị hiện hành cho ứng dụng web.

\section{Giải pháp xác thực tập trung: Keycloak}
\label{sec:keycloak}

\subsection{Giới thiệu Keycloak}

Keycloak~\cite{keycloak} là một giải pháp mã nguồn mở về quản lý định danh và
truy cập, hỗ trợ sẵn SSO, OIDC, OAuth 2.0, SAML 2.0, xác thực đa yếu tố và
liên kết với các nguồn định danh bên ngoài (LDAP, Active Directory, mạng xã
hội). Keycloak được viết trên nền Quarkus, có khả năng mở rộng và được sử
dụng rộng rãi trong cả môi trường doanh nghiệp lẫn học thuật. Đây là lý do
Keycloak được lựa chọn làm nền tảng IdP cho chuyên đề.

\subsection{Các khái niệm chính}

\begin{itemize}
  \item \textbf{Realm:} một không gian quản lý định danh độc lập, chứa người
  dùng, vai trò và các client.
  \item \textbf{Client:} một ứng dụng hoặc dịch vụ được phép sử dụng Keycloak
  để xác thực.
  \item \textbf{Role:} vai trò dùng cho ủy quyền (RBAC).
  \item \textbf{User:} tài khoản người dùng thuộc một realm.
  \item \textbf{Admin REST API:} giao diện lập trình cho phép tự động hóa cấu
  hình và quản trị Keycloak.
\end{itemize}

\subsection{So sánh với một số giải pháp khác}

Bảng~\ref{tab:ssosanh-idp} so sánh Keycloak với một số giải pháp định danh
phổ biến.

\begin{table}[htbp]
\centering
\caption{So sánh Keycloak với một số giải pháp định danh}
\label{tab:ssosanh-idp}
\small
\begin{tabular}{p{2.6cm} c c c c}
\toprule
\textbf{Tiêu chí} & \textbf{Keycloak} & \textbf{Okta/Auth0} & \textbf{Ory Hydra} & \textbf{Authentik}\\
\midrule
Mã nguồn mở     & Có (Apache) & Không & Có & Có\\
Tự vận hành     & Có     & Hạn chế & Có    & Có\\
OIDC/OAuth2/SAML& Có     & Có      & OAuth2/OIDC & Có\\
Giao diện quản trị & Có  & Có      & Tối giản & Có\\
Độ phổ biến / tài liệu & Cao & Cao & Trung bình & Trung bình\\
Chi phí         & Miễn phí & Trả phí & Miễn phí & Miễn phí\\
\bottomrule
\end{tabular}
\end{table}

\subsection{Kiến trúc nội bộ của Keycloak}
\label{subsec:kc-arch}

Keycloak được xây dựng trên nền tảng Quarkus và tổ chức thành các lớp: lớp điểm
cuối REST xử lý yêu cầu giao thức, lớp dịch vụ chứa logic nghiệp vụ, lớp lưu
trữ (JPA) làm việc với cơ sở dữ liệu, và lớp bộ nhớ đệm Infinispan dùng cho
phiên đăng nhập cùng dữ liệu realm. Kiến trúc này cho phép Keycloak chạy nhiều
bản sao trong mô hình cụm để đảm bảo tính sẵn sàng cao.

Keycloak cung cấp nhiều điểm mở rộng (SPI) cho xác thực, ánh xạ giao thức và
lưu trữ, cùng Admin REST API cho phép tự động hóa toàn bộ thao tác quản trị.
Nó cũng hỗ trợ xuất/nhập cấu hình realm dưới dạng JSON và cung cấp các điểm
cuối về sức khỏe (health) và chỉ số (metrics) phục vụ giám sát. Những đặc điểm
này là nền tảng để quản lý cấu hình như mã nguồn và triển khai tự động.

\subsection{Các chuẩn và giao thức liên quan}
\label{subsec:related}

Ngoài OAuth 2.0 và OIDC, trong lĩnh vực định danh còn tồn tại nhiều chuẩn khác:

\begin{itemize}
  \item \textit{SAML 2.0}: giao thức định danh liên kết dựa trên XML, phổ biến
  trong môi trường doanh nghiệp.
  \item \textit{LDAP}: giao thức truy cập thư mục, thường được dùng làm nguồn
  người dùng.
  \item \textit{Kerberos}: giao thức xác thực phổ biến trong mạng nội bộ.
  \item \textit{WebAuthn}: chuẩn xác thực không mật khẩu dựa trên khóa công
  khai.
\end{itemize}

Keycloak hỗ trợ liên kết với LDAP/Active Directory và SAML, nhờ đó có thể tích
hợp vào hạ tầng định danh sẵn có của tổ chức mà không cần thay thế toàn bộ.

\section{Ảo hóa container và Kubernetes}
\label{sec:k8s}

\subsection{Container hóa với Docker}

Docker đóng gói ứng dụng cùng các phụ thuộc vào một \textit{container} -- một
đơn vị chạy nhẹ, nhất quán giữa các môi trường. Image của Keycloak, PostgreSQL
và ứng dụng demo đều có thể được phân phối dưới dạng image, giúp việc triển
khai trở nên tái lập được.

\subsection{Kubernetes và k3s}

Kubernetes~\cite{k8s} là nền tảng điều phối container, cung cấp khả năng tự
phục hồi, mở rộng và quản lý vòng đời dịch vụ. Trong phạm vi chuyên đề, bản
phân phối nhẹ k3s được sử dụng để dựng cụm Kubernetes ngay trên máy
cá nhân hoặc trên máy ảo, phù hợp với điều kiện thực hành hạn chế về hạ tầng.

\subsection{Helm}

Helm là trình quản lý gói cho Kubernetes, cho phép đóng gói các manifest
thành \textit{chart} có thể tham số hóa và tái sử dụng. Keycloak có chart chính
thức, thuận lợi cho việc triển khai tự động.

\subsection{Các đối tượng cốt lõi của Kubernetes}
\label{subsec:k8s-obj}

Kubernetes quản lý ứng dụng thông qua các đối tượng khai báo. Một số đối tượng
quan trọng:

\begin{itemize}
  \item \textit{Pod}: đơn vị nhỏ nhất, chứa một hoặc nhiều container cùng chạy.
  \item \textit{Deployment}: quản lý việc tạo và cập nhật các bản sao của Pod,
  hỗ trợ cuốn chiếu và tự phục hồi.
  \item \textit{Service}: cung cấp địa chỉ ổn định và cân bằng tải cho một nhóm
  Pod.
  \item \textit{Ingress}: định tuyến lưu lượng HTTP/HTTPS từ bên ngoài vào các
  Service.
  \item \textit{ConfigMap} và \textit{Secret}: tách cấu hình và thông tin nhạy
  cảm khỏi image.
  \item \textit{StatefulSet} và \textit{PersistentVolume}: quản lý các ứng dụng
  có trạng thái như cơ sở dữ liệu.
  \item \textit{HorizontalPodAutoscaler}: tự động mở rộng số bản sao theo tải.
\end{itemize}

Việc hiểu rõ các đối tượng này là điều kiện cần để xây dựng manifest cho dịch
vụ định danh và ứng dụng demo.

\subsection{Container hoạt động như thế nào}
\label{subsec:container}

Container là một tiến trình được cô lập nhờ các tính năng của nhân Linux:
\textit{namespace} (cô lập không gian tên, mạng, tiến trình) và \textit{cgroup}
(giới hạn tài nguyên CPU, bộ nhớ). Khác với máy ảo, container chia sẻ nhân với
máy chủ nên nhẹ và khởi động nhanh. Image được tổ chức thành nhiều lớp (layer)
xếp chồng, giúp tái sử dụng và phân phối hiệu quả.

\subsection{Kiến trúc của Kubernetes}
\label{subsec:k8sarch}

Một cụm Kubernetes gồm hai phần: \textit{control plane} (máy chủ điều khiển) và
các \textit{node} (máy công tác). Control plane chứa API server (điểm vào duy
nhất), etcd (lưu trạng thái), scheduler (phân bổ Pod) và controller manager
(duy trì trạng thái mong muốn). Trên mỗi node, kubelet quản lý container và
kube-proxy xử lý mạng. k3s là bản phân phối gọn nhẹ, đóng gói các thành phần
này vào một tiến trình duy nhất, phù hợp cho môi trường học tập và thiết bị hạn
chế.

\section{Triển khai tự động}
\label{sec:autodeploy}

\subsection{Tích hợp và triển khai liên tục (CI/CD)}

CI/CD tự động hóa việc kiểm thử, đóng gói và phát hành phần mềm. Mỗi thay đổi
được đẩy lên kho mã nguồn sẽ kích hoạt một \textit{pipeline} thực hiện: kiểm tra
mã (lint), chạy kiểm thử, xây dựng image, quét lỗ hổng và phát hành. Điều này
giúp giảm sai sót do thao tác thủ công.

\subsection{Hạ tầng dưới dạng mã (IaC)}

IaC (Infrastructure as Code) mô tả hạ tầng bằng các tệp cấu hình (ví dụ
Terraform, Ansible) thay vì thao tác thủ công, cho phép tái lập môi trường và
quản lý hạ tầng bằng quy trình phiên bản hóa.

\subsection{GitOps và ArgoCD}

GitOps lấy kho Git làm nguồn chân lý duy nhất cho trạng thái mong muốn
của hệ thống. ArgoCD~\cite{argocd} liên tục so sánh trạng thái thực tế trên
cụm Kubernetes với trạng thái trong Git và tự động đồng bộ, giúp việc triển
khai trở nên minh bạch, có thể truy vết và dễ dàng khôi phục.

\subsection{Quản lý cấu hình như mã nguồn}

Đối với Keycloak, cấu hình định danh (realm, client, role, chính sách) cũng có
thể được quản lý như mã nguồn thông qua các công cụ như
\texttt{keycloak-config-cli}~\cite{kc-config-cli} hoặc Keycloak Terraform
Provider~\cite{tf-keycloak}. Đây là cầu nối giữa dịch vụ định danh và quy
trình triển khai tự động, là trọng tâm của chuyên đề.

\subsection{Lợi ích và thách thức của triển khai tự động}
\label{subsec:auto-tradeoff}

Triển khai tự động mang lại nhiều lợi ích: khả năng tái lập môi trường, tính
nhất quán giữa các lần triển khai, giảm thời gian khắc phục sự cố và tăng tốc
độ phát hành. Mọi thay đổi đều có thể truy vết qua lịch sử Git, giúp việc kiểm
toán và khôi phục trở nên dễ dàng.

Tuy nhiên, hướng tiếp cận này cũng đặt ra thách thức: đòi hỏi kỹ năng và thời
gian đầu tư ban đầu, việc quản lý bí mật (khóa, mật khẩu) trở nên phức tạp hơn,
bản thân pipeline cũng cần được kiểm thử, và lỗi cấu hình có thể lan rộng
nhanh. Do đó, nguyên tắc ``tự động hóa có kiểm soát'' -- kết hợp kiểm thử tự
động và khả năng quay lui -- là cần thiết.

\subsection{Ví dụ một quy trình CI/CD điển hình}
\label{subsec:pipeline}

Một quy trình CI/CD hoàn chỉnh thường trải qua các giai đoạn sau: (1) lấy mã
nguồn và kiểm tra định dạng (lint); (2) chạy kiểm thử đơn vị; (3) xây dựng
image; (4) quét lỗ hổng bảo mật; (5) đẩy image lên registry; (6) cập nhật
manifest và để GitOps đồng bộ; (7) kiểm tra sau triển khai (smoke test). Mỗi
giai đoạn có thể chặn quy trình nếu thất bại, đảm bảo chỉ những thay đổi đạt
chất lượng mới được phát hành.

\section{Kết chương}

Chương này đã trình bày cơ sở lý thuyết về xác thực tập trung, giao thức
OAuth 2.0/OIDC, nền tảng Keycloak và các phương pháp triển khai tự động. Trên
nền tảng đó, chương tiếp theo sẽ phân tích yêu cầu và thiết kế kiến trúc cho
hệ thống cụ thể của chuyên đề.
